\documentclass[aspectratio=169]{beamer}
\usetheme{Warsaw}
\usepackage[utf8]{inputenc}
\usepackage[T1]{fontenc}
\usepackage{amsmath, amssymb}
\usepackage{booktabs}
\usepackage{graphicx}

% ── Colours ───────────────────────────────────────────────────────────────────
\setbeamercolor{structure}{fg=blue!70!black}

% ── Metadata ──────────────────────────────────────────────────────────────────
\title{Academic Seminar Presentation}
\subtitle{Warsaw Theme for Formal Settings}
\author{Prof.\ Maria García}
\institute{Department of Computer Science \\ University of Example}
\date{Academic Year 2025--2026}

\begin{document}

\begin{frame}
  \titlepage
\end{frame}

\begin{frame}{Outline}
  \tableofcontents
\end{frame}

% ── Section 1 ─────────────────────────────────────────────────────────────────
\section{Background \& Related Work}

\begin{frame}{Foundational Concepts}
  \begin{definition}
    A \textbf{system} $\mathcal{S}$ is defined as a tuple $(Q, \Sigma, \delta, q_0, F)$ where:
    \begin{itemize}
      \item $Q$ — finite set of states
      \item $\Sigma$ — input alphabet
      \item $\delta : Q \times \Sigma \to Q$ — transition function
    \end{itemize}
  \end{definition}
\end{frame}

\begin{frame}{Related Work}
  \begin{block}{Classical approaches}
    Early solutions relied on heuristic methods~\cite{example2020}, achieving moderate accuracy
    but suffering from scalability issues.
  \end{block}

  \begin{block}{Recent advances}
    Deep learning methods have significantly improved performance, with~\cite{modern2024}
    reporting a 15\% gain over prior work.
  \end{block}
\end{frame}

% ── Section 2 ─────────────────────────────────────────────────────────────────
\section{Theoretical Framework}

\begin{frame}{Main Theorem}
  \begin{theorem}[Convergence]
    Let $\{x_k\}$ be a sequence generated by Algorithm~1. Under assumptions A1--A3,
    \[
      \|x_k - x^*\| \leq C \cdot \rho^k, \quad \rho \in (0,1),
    \]
    where $x^*$ is the optimal solution and $C > 0$ is a constant.
  \end{theorem}

  \begin{proof}[Sketch]
    By Lemma~2 and the contraction mapping principle, the sequence converges
    geometrically to the fixed point $x^*$.  \qed
  \end{proof}
\end{frame}

\begin{frame}{Algorithm}
  \begin{columns}
    \column{0.55\textwidth}
    \textbf{Algorithm 1: Iterative Solver}
    \begin{enumerate}
      \item Initialise $x_0 \in \mathcal{X}$, $k \leftarrow 0$
      \item \textbf{While} $\|x_k - x_{k-1}\| > \varepsilon$:
        \begin{enumerate}
          \item Compute gradient $\nabla f(x_k)$
          \item Update $x_{k+1} \leftarrow x_k - \alpha_k \nabla f(x_k)$
          \item $k \leftarrow k+1$
        \end{enumerate}
      \item \textbf{Return} $x_k$
    \end{enumerate}

    \column{0.4\textwidth}
    \textbf{Complexity:}
    \begin{itemize}
      \item Time: $\mathcal{O}(n^2 \log n)$
      \item Space: $\mathcal{O}(n)$
    \end{itemize}
  \end{columns}
\end{frame}

% ── Section 3 ─────────────────────────────────────────────────────────────────
\section{Experimental Evaluation}

\begin{frame}{Experimental Setup}
  \begin{itemize}
    \item \textbf{Datasets:} Benchmark A (10k instances), Benchmark B (50k instances).
    \item \textbf{Baselines:} Method X, Method Y, State-of-the-art Z.
    \item \textbf{Metrics:} Accuracy, F1-score, runtime.
    \item \textbf{Hardware:} 4$\times$ NVIDIA A100 GPUs, 512\,GB RAM.
  \end{itemize}
\end{frame}

\begin{frame}{Results}
  \begin{table}
    \centering
    \begin{tabular}{lccc}
      \toprule
      \textbf{Method} & \textbf{Acc.} & \textbf{F1} & \textbf{Time (s)} \\
      \midrule
      Method X        & 78.4          & 0.76        & 12.3              \\
      Method Y        & 82.1          & 0.81        & 18.7              \\
      SOTA Z          & 85.9          & 0.84        & 45.2              \\
      \textbf{Ours}   & \textbf{89.3} & \textbf{0.88} & \textbf{9.8}  \\
      \bottomrule
    \end{tabular}
    \caption{Performance on Benchmark A.}
  \end{table}
\end{frame}

% ── Conclusion ────────────────────────────────────────────────────────────────
\section{Conclusions}

\begin{frame}{Conclusions \& Future Work}
  \textbf{We have shown:}
  \begin{itemize}
    \item Our approach achieves SOTA with 3.4$\times$ lower runtime.
    \item The theoretical guarantee (Theorem~1) holds empirically.
  \end{itemize}

  \vspace{0.8em}
  \textbf{Future directions:}
  \begin{itemize}
    \item Extension to online/streaming settings.
    \item Application to domain-specific datasets.
  \end{itemize}
\end{frame}

\begin{frame}{References}
  \small
  \begin{thebibliography}{9}
    \bibitem{example2020} A. Author, ``Classical Heuristics,'' \textit{ICML}, 2020.
    \bibitem{modern2024} B. Author, ``Deep Methods,'' \textit{NeurIPS}, 2024.
  \end{thebibliography}
\end{frame}

\begin{frame}[plain]
  \centering
  \vfill
  {\Huge\bfseries Thank You} \\[1em]
  {\large\itshape Questions \& Discussion}
  \vfill
\end{frame}

\end{document}
